% !PRINTED 601
\weberpage{601}
Let us take in a plane a rectangular coordinate system
 $X, Y$, and let us call \gesperrt{lattice points} those points
with the coordinates
\[
X = (x_0 + Mx)\, a_{22}\, \sqrt{\dfrac{a}{t}},
\]
\[
Y = (y_0 + My)\, a_{22}\, \sqrt{\dfrac{a}{t}},
\]
where $x_0, y_0$ are kept fixed and $x, y$ run through all integral
rational numbers; then by this lattice the plane is divided into
squares whose area is
\[
\delta^2 = M^2\, \dfrac{a_{22}^2\, a}{t}
\]
and if $T_0$ is the number of lattice points that lie within
the ellipse
\[
\varphi(X, Y) = 1
\]
or on its boundary, then $T_0\,\delta^2$ is, for an infinite $t$,
equal to the area of this ellipse, and hence equal to $\dfrac{\pi}{\sqrt{-\Delta}}$.

If one takes account further of the elements of area $\delta^2$ that are
cut through by the periphery of the ellipse, from which one obtains a
measure for the error of this determination of area, then
\weq{5}{T_0 = \dfrac{\pi t}{M^2\, a_{22}^2\, a\, \sqrt{-\Delta}} + R_0\sqrt{t},}
where $R_0$ is a quantity that remains finite as $t$ increases
infinitely [Bd.\,II, \S\,194, (6)].

We obtain such an expression for every $\alpha_0$, and if
we add all these expressions together, we obtain, having regard
to (3), the theorem:

\gesperrt{If $A$ is a congruence group with the modulus $M$
and $\mathfrak{m}$ an ideal in $\Omega$ coprime to $M$, and further $T$
the number of integral principal ideals contained in $\overline{E}A$
whose norm is not greater than $t$, then}
\weq{6}{T = \dfrac{g t}{N(\mathfrak{m})} + R\sqrt{t},}
\gesperrt{where $g$ is a number independent of $\mathfrak{m}$ and $t$ and different
from zero, and $R$ is a function of $t$ that is finite as $A$ % [sic]
increases infinitely.}
